\chapter{เทอร์โมไดนามิกส์การเผาไหม้และฟิสิกส์เชื้อเพลิงชีวมวล}
\label{chap:ch05}

\section*{แผนบริหารการสอนประจำบทที่ 5}
\addcontentsline{toc}{section}{แผนบริหารการสอนประจำบทที่ 5}

\noindent\textbf{หัวข้อเนื้อหาประจำบท}
\begin{enumerate}
    \item กฎทางอุณหพลศาสตร์และสมดุลพลังงานการเผาไหม้
    \item การวิเคราะห์คุณสมบัติเชื้อเพลิงชีวมวลโดยประมาณ
    \item ค่าความร้อนของเชื้อเพลิงและการวัดด้วยบอมบ์แคลอริมิเตอร์
    \item ชีวกลศาสตร์ของการอัดเม็ดและอัดแท่งชีวมวล
\end{enumerate}

\noindent\textbf{วัตถุประสงค์เชิงพฤติกรรม}
\begin{enumerate}
    \item อธิบายหลักการ ทฤษฎีฟิสิกส์ และแบบจำลองทางสิ่งแวดล้อมประจำบทที่ 5 ได้อย่างถูกต้อง
    \item วิเคราะห์ สังเคราะห์ และประยุกต์ใช้เครื่องมือตรวจวัดหรือกระบวนการแปรรูปพลังงานได้อย่างมีประสิทธิภาพ
    \item คำนวณ ประเมินผลทางเทอร์โมไดนามิกส์ ทัศนศาสตร์ หรือการลดก๊าซเรือนกระจกได้อย่างถูกต้องตามหลักวิชาการ
\end{enumerate}

\noindent\textbf{การวัดและประเมินผล}
\begin{table}[H]
\centering\small
\begin{tabularx}{\textwidth}{p{0.55\textwidth} p{0.25\textwidth} c}
\toprule
\textbf{เกณฑ์การประเมินผล} & \textbf{เครื่องมือวัดผล} & \textbf{สัดส่วนคะแนน} \\
\midrule
1. ทักษะการปฏิบัติการทดลองและการคำนวณ & แบบประเมินทักษะห้องแล็บ & 40\% \\
2. รายงานผลการทดลองและการวิเคราะห์ & การตรวจดาต้าชีทและรายงาน & 30\% \\
3. แบบฝึกหัดทบทวนและคำถามท้ายบท & แบบทดสอบท้ายบทเรียน & 20\% \\
4. จิตพิสัยและการมีส่วนร่วม & การเข้าเรียนและการตรงต่อเวลา & 10\% \\
\midrule
\textbf{รวมทั้งสิ้น} & & \textbf{100\%} \\
\bottomrule
\end{tabularx}
\end{table}
\newpage

% ==========================================================
% เนื้อหาบทเรียนประจำบทที่ 5
% ==========================================================

\section{กฎทางอุณหพลศาสตร์และสมดุลพลังงานการเผาไหม้}
\index{กฎทางอุณหพลศาสตร์และสมดุลพลังงานการเผาไหม้}

การศึกษาและทำความเข้าใจเกี่ยวกับกฎทางอุณหพลศาสตร์และสมดุลพลังงานการเผาไหม้ (Combustion Thermodynamics and Energy Balance) มีความสำคัญอย่างยิ่งในงานฟิสิกส์สิ่งแวดล้อมและพลังงาน โดยมีรายละเอียดสาระสำคัญดังนี้

การเผาไหม้ (Combustion) เป็นปฏิกิริยาเคมีคายความร้อนอย่างรวดเร็วระหว่างเชื้อเพลิงที่มีคาร์บอนและไฮโดรเจนเป็นองค์ประกอบหลักกับสารออกซิไดซ์ (ออกซิเจนในอากาศ)

สมการเคมีพื้นฐานของการเผาไหม้สมบูรณ์:
\[ \text{C}\_x\text{H}\_y\text{O}\_z + \left(x + \frac{y}{4} - \frac{z}{2}\right) \text{O}\_2 \rightarrow x\text{CO}\_2 + \frac{y}{2}\text{H}\_2\text{O} + \Delta H\_c \]
โดยที่ $\Delta H_c$ คือเอนทาลปีของการเผาไหม้ (Enthalpy of combustion)

สมดุลความร้อนของเตาเผาชีวมวลเป็นไปตามกฎข้อที่หนึ่งของเทอร์โมไดนามิกส์: พลังงานความร้อนที่ปลดปล่อยออกมาจะถูกนำไปใช้เพิ่มอุณหภูมิของผลิตภัณฑ์ก๊าซไอเสียและถ่ายเทความร้อนสู่ระบบแลกเปลี่ยนความร้อน การเผาไหม้ที่ไม่สมบูรณ์อันเนื่องมาจากออกซิเจนไม่เพียงพอหรืออุณหภูมิในห้องเผาไหม้ต่ำเกินไป จะก่อให้เกิดก๊าซคาร์บอนมอนอกไซด์ (CO), ไฮโดรคาร์บอนที่เผาไหม้ไม่หมด (UHC), และอนุภาคเขม่าควันดำ (Soot)

\begin{physbox}[หลักการและสมการฟิสิกส์สิ่งแวดล้อม]
การทำความเข้าใจพฤติกรรมของอนุภาคและการถ่ายเทความร้อนในระบบสิ่งแวดล้อมต้องอาศัยการประยุกต์ใช้กฎเทอร์โมไดนามิกส์และทัศนศาสตร์คลื่นแม่เหล็กไฟฟ้าอย่างแม่นยำ
\end{physbox}

\section{การวิเคราะห์คุณสมบัติเชื้อเพลิงชีวมวลโดยประมาณ}
\index{การวิเคราะห์คุณสมบัติเชื้อเพลิงชีวมวลโดยประมาณ}

การศึกษาและทำความเข้าใจเกี่ยวกับการวิเคราะห์คุณสมบัติเชื้อเพลิงชีวมวลโดยประมาณ (Proximate and Ultimate Analysis of Biomass) มีความสำคัญอย่างยิ่งในงานฟิสิกส์สิ่งแวดล้อมและพลังงาน โดยมีรายละเอียดสาระสำคัญดังนี้

การประเมินคุณภาพเชื้อเพลิงชีวมวลแข็งอาศัยการวิเคราะห์หลัก 2 รูปแบบ:

1. การวิเคราะห์โดยประมาณ (Proximate Analysis):
- ปริมาณความชื้น (Moisture Content; MC): ส่งผลกระทบเชิงลบต่อค่าความร้อน เนื่องจากต้องสูญเสียพลังงานความร้อนแฝงในการระเหยน้ำ
- สารระเหยง่าย (Volatile Matter; VM): ก๊าซที่ระเหยออกมาระหว่างการให้ความร้อนที่ 900$^{\circ}$C โดยไม่มีอากาศ ชีวมวลส่วนใหญ่มี VM สูง (70-80\%) ทำให้ติดไฟง่ายแต่เกิดเปลวไฟยาวและมีควัน
- คาร์บอนคงตัว (Fixed Carbon; FC): ส่วนของคาร์บอนแข็งที่เหลือหลังการระเหยสารระเหย เป็นตัวให้พลังงานความร้อนหลักและมอดช้า
- ปริมาณเถ้า (Ash Content): สารอนินทรีย์ที่ไม่เผาไหม้ เช่น ซิลิกาและเหล็กออกไซด์

2. การวิเคราะห์ธาตุองค์ประกอบ (Ultimate Analysis): วิเคราะห์สัดส่วนร้อยละของธาตุ C, H, O, N, และ S

\begin{flamebox}[นวัตกรรมพลังงานความร้อน มรภ.รำไพพรรณี]
การเสริมประสิทธิภาพความร้อนของถ่านอัดแท่งด้วยคาร์บอนแบล็คช่วยยกระดับค่าความร้อนสูงเกิน 6,200 cal/g ลดควันดำ และเพิ่มความแข็งแรงเชิงกล ตอบโจทย์การใช้งานเชิงพาณิชย์และสิ่งแวดล้อม
\end{flamebox}

\section{ค่าความร้อนของเชื้อเพลิงและการวัดด้วยบอมบ์แคลอริมิเตอร์}
\index{ค่าความร้อนของเชื้อเพลิงและการวัดด้วยบอมบ์แคลอริมิเตอร์}

การศึกษาและทำความเข้าใจเกี่ยวกับค่าความร้อนของเชื้อเพลิงและการวัดด้วยบอมบ์แคลอริมิเตอร์ (Heating Values and Bomb Calorimetry) มีความสำคัญอย่างยิ่งในงานฟิสิกส์สิ่งแวดล้อมและพลังงาน โดยมีรายละเอียดสาระสำคัญดังนี้

พลังงานความร้อนของเชื้อเพลิงแบ่งออกเป็น:
- ค่าความร้อนสูง (Higher Heating Value; HHV หรือ Gross Calorific Value; GCV): ค่าพลังงานความร้อนทั้งหมดที่ได้จากการเผาไหม้ โดยไอน้ำที่เกิดขึ้นควบแน่นกลับเป็นน้ำเหลวและปลดปล่อยความร้อนแฝงคืนสู่ระบบ
- ค่าความร้อนต่ำ (Lower Heating Value; LHV หรือ Net Calorific Value; NCV): ค่าพลังงานความร้อนจริงในการใช้งานเมื่อไอน้ำระบายออกไปในรูปก๊าซ:
\[ \text{LHV} = \text{HHV} - 2.442 \times (9H + \text{MC}) \]
การวัดค่า HHV ในห้องปฏิบัติการทำโดยใช้เครื่อง ออกซิเจนบอมบ์แคลอริมิเตอร์ (Oxygen Bomb Calorimeter) ภายใต้บรรยากาศออกซิเจนบริสุทธิ์ความดัน 30 บรรยากาศ

\section{ชีวกลศาสตร์ของการอัดเม็ดและอัดแท่งชีวมวล}
\index{ชีวกลศาสตร์ของการอัดเม็ดและอัดแท่งชีวมวล}

การศึกษาและทำความเข้าใจเกี่ยวกับชีวกลศาสตร์ของการอัดเม็ดและอัดแท่งชีวมวล (Densification and Briquetting Mechanics) มีความสำคัญอย่างยิ่งในงานฟิสิกส์สิ่งแวดล้อมและพลังงาน โดยมีรายละเอียดสาระสำคัญดังนี้

ชีวมวลทางการเกษตรมีความหนาแน่นรวมต่ำ (Bulk density $< 200$ kg/m$^3$) ทำให้การขนส่งและการเก็บรักษาสิ้นเปลืองค่าใช้จ่าย กระบวนการอัดแท่ง (Briquetting) อาศัยแรงอัดทางกลสูง (Mechanical pressure) และความร้อนเพื่อหลอมละลายสารลิกนิน (Lignin) ตามธรรมชาติในชีวมวลให้ทำหน้าที่เป็นตัวประสาน (Natural binder) เพิ่มความหนาแน่นรวมเป็นมากกว่า 1,000 kg/m$^3$

\begin{flamebox}[นวัตกรรมพลังงานความร้อน มรภ.รำไพพรรณี]
การเสริมประสิทธิภาพความร้อนของถ่านอัดแท่งด้วยคาร์บอนแบล็คช่วยยกระดับค่าความร้อนสูงเกิน 6,200 cal/g ลดควันดำ และเพิ่มความแข็งแรงเชิงกล ตอบโจทย์การใช้งานเชิงพาณิชย์และสิ่งแวดล้อม
\end{flamebox}

\section{ขั้นตอนและวิธีปฏิบัติการทดลอง}
\index{ขั้นตอนการทดลองประจำบทที่ 5}

\subsection{การวิเคราะห์ปริมาณความชื้นและเถ้าของเชื้อเพลิงชีวมวลตามมาตรฐาน ASTM}
\begin{conceptbox}[การวิเคราะห์ปริมาณความชื้นและเถ้าของเชื้อเพลิงชีวมวลตามมาตรฐาน ASTM]
1. ชั่งตัวอย่างชีวมวลบดละเอียด 1.0000 กรัม ใส่ในถ้วยหาความชื้นที่ทราบน้ำหนัก
2. อบในตู้อบลมร้อนที่ 105$^{\circ}$C เป็นเวลา 2 ชั่วโมง ปล่อยให้เย็นในโถดูดความชื้นแล้วชั่งน้ำหนัก คำนวณร้อยละความชื้น
3. นำตัวอย่างแห้งไปเผาในเตาเผาอุณหภูมิสูง (Muffle furnace) ที่ 750$^{\circ}$C เป็นเวลา 4 ชั่วโมงจนเถ้าเป็นสีขาวเทา
4. ปล่อยให้เย็นและชั่งน้ำหนักเถ้าที่เหลือ คำนวณร้อยละของปริมาณเถ้า (\% Ash)
\end{conceptbox}

\subsection{การทดสอบค่าความร้อนสูง ด้วยเครื่อง Oxygen Bomb Calorimeter}
\begin{conceptbox}[การทดสอบค่าความร้อนสูง ด้วยเครื่อง Oxygen Bomb Calorimeter]
1. ชั่งตัวอย่างชีวมวลอัดเม็ด 1.0 กรัม ผูกลวดฟิวส์นิกเกิล-โครเมียม สัมผัสตัวอย่างในเบ้าหลอม
2. ประกอบบอมบ์โลหะสเตนเลส อัดก๊าซออกซิเจนบริสุทธิ์เข้าสู่บอมบ์จนได้ความดัน 3.0 MPa (30 bar)
3. จุ่มบอมบ์ลงในถังน้ำแคลอริมิเตอร์ที่มีน้ำกลั่น 2,000 กรัม ควบคุมอุณหภูมิเริ่มต้น
4. จุดชนวนด้วยกระแสไฟฟ้า บันทึกการเพิ่มขึ้นของอุณหภูมิน้ำ ($\Delta T$) ละเอียดระดับ 0.001$^{\circ}$C
5. คำนวณค่า HHV (cal/g หรือ MJ/kg) โดยชดเชยค่าความจุความร้อนของเครื่องและพลังงานจากลวดจุดชนวน
\end{conceptbox}

\section{ตารางบันทึกผลการทดลองและการอภิปรายผล}
\begin{table}[H]
\centering\small
\caption{ตารางบันทึกผลการทดลองและการตรวจวัดพารามิเตอร์ประจำบทที่ 5}
\label{tab:ch05_datasheet}
\begin{tabularx}{\textwidth}{p{0.3\textwidth} p{0.3\textwidth} X}
\toprule
\textbf{พารามิเตอร์ที่ตรวจวัด} & \textbf{ค่าที่บันทึกได้} & \textbf{การแปลผลและการวิเคราะห์ทางฟิสิกส์} \\
\midrule
การทดสอบชุดที่ 1 & ค่าอยู่ในเกณฑ์มาตรฐาน & สอดคล้องกับแบบจำลองทฤษฎี \\
การทดสอบชุดที่ 2 & ค่ามีแนวโน้มเพิ่มขึ้นตามสัดส่วน & ยืนยันสมมติฐานการวิจัยเชิงประจักษ์ \\
ประสิทธิภาพรวม & ร้อยละ 88.5 & แสดงผลสัมฤทธิ์ในระดับยอมรับได้ \\
\bottomrule
\end{tabularx}
\end{table}

\section{คำถามทบทวนและแบบฝึกหัดท้ายบท}
\begin{enumerate}
    \item จงอธิบายความแตกต่างระหว่างค่าความร้อนสูง (HHV) และค่าความร้อนต่ำ (LHV) ของเชื้อเพลิงชีวมวล
    \item เหตุใดสารระเหยง่าย (Volatile matter) ที่สูงในชีวมวลจึงส่งผลให้เกิดเขม่าควันได้ง่ายกว่าถ่านหินแอนทราไซต์
    \item จงคำนวณค่าความร้อนต่ำ (LHV) หากเชื้อเพลิงชีวมวลมีค่า $\text{HHV} = 4,200$ cal/g, ปริมาณไฮโดรเจน $H = 5.5\%$, และความชื้น $\text{MC} = 10\%$
\end{enumerate}
